\documentclass[12pt,a4paper]{article}
% Preâmbulo institucional comum — DEE/IFMA
\usepackage[utf8]{inputenc}
\usepackage[T1]{fontenc}
\usepackage[brazil]{babel}
\usepackage{lmodern}
\usepackage{microtype}
\usepackage{geometry}
\usepackage{xcolor}
\usepackage{graphicx}
\usepackage{booktabs}
\usepackage{tabularx}
\usepackage{array}
\usepackage{enumitem}
\usepackage{hyperref}
\usepackage{lastpage}
\usepackage{fancyhdr}
\usepackage{setspace}
\usepackage{titlesec}
\usepackage{tcolorbox}
\usepackage{ragged2e}

\geometry{a4paper,top=2.2cm,bottom=2.2cm,left=2.5cm,right=2.5cm}
\definecolor{azulpetroleo}{HTML}{0F4C5C}
\definecolor{ambar}{HTML}{C8892D}
\definecolor{cinzaclaro}{HTML}{F2F5F7}
\definecolor{cinzaescuro}{HTML}{374151}
\hypersetup{colorlinks=true,linkcolor=azulpetroleo,urlcolor=azulpetroleo,citecolor=azulpetroleo}
\setlength{\parindent}{1.25cm}
\setlength{\parskip}{0.35em}
\onehalfspacing
\titleformat{\section}{\large\bfseries\color{azulpetroleo}}{}{0pt}{}
\titleformat{\subsection}{\normalsize\bfseries\color{azulpetroleo}}{}{0pt}{}
\pagestyle{fancy}
\fancyhf{}
\fancyhead[L]{\small Departamento de Eletroeletrônica — DEE}
\fancyhead[R]{\small IFMA — Campus São Luís/Monte Castelo}
\fancyfoot[C]{\small Página \thepage\ de \pageref{LastPage}}
\renewcommand{\headrulewidth}{0.4pt}
\renewcommand{\footrulewidth}{0.4pt}
\newcommand{\campo}[1]{\textcolor{ambar}{\textbf{[#1]}}}
\newcommand{\instituicao}{Instituto Federal de Educação, Ciência e Tecnologia do Maranhão}
\newcommand{\campus}{Campus São Luís/Monte Castelo}
\newcommand{\departamento}{Departamento de Eletroeletrônica — DEE}
\newcommand{\cidade}{São Luís — MA}
\newcommand{\cabecalhoInstitucional}{%
\begin{center}
{\Large\bfseries\color{azulpetroleo}\instituicao}\\[2mm]
{\large\bfseries\campus}\\[1mm]
{\normalsize\departamento}
\end{center}
\vspace{2mm}
\hrule
\vspace{4mm}
}
\newtcolorbox{destaque}{colback=cinzaclaro,colframe=azulpetroleo,boxrule=0.6pt,arc=2mm}

\newcommand{\empresa}{Texas Instruments}
\newcommand{\prioridade}{P1 — Máxima}
\newcommand{\setorContato}{University Program / Power Management / Embedded Processing}
\newcommand{\areaEmpresa}{eletrônica de potência, conversores, controle digital, microcontroladores, DSPs, instrumentação e sistemas embarcados}
\newcommand{\estrategia}{Prospecção educacional prioritária com base no portfólio universitário de Eletrônica de Potência e plataformas de controle digital.}
\newcommand{\siteOficial}{\url{https://www.ti.com/}}
\newcommand{\textoCooperacao}{A Texas Instruments mantém conteúdo e kits voltados ao ensino universitário de Eletrônica de Potência, incluindo experimentos com conversores e controle. Propõe-se estruturar bancadas para conversão CC--CC, controle digital, acionamentos, fontes chaveadas e sistemas embarcados.}
\newcommand{\itensEmpresa}{Kits TI-PMLK Buck/Boost/BuckBoost ou equivalentes & 6 & Conversores CC--CC e fundamentos de eletrônica de potência\\
Placas C2000 LaunchPad/controlCARD ou equivalentes & 8 & Controle digital de conversores e motores\\
Kits/placas de GaN e drivers de potência & 6 & Conversores de alta frequência e eficiência\\
Placas de avaliação de fontes e reguladores & 8 & Fontes chaveadas e gerenciamento de energia\\
Kits de aquisição e sensoriamento de corrente/tensão & 1 lote & Instrumentação de conversores e acionamentos\\
Componentes e amostras para prototipagem & 1 lote & Projetos, TCCs e pesquisa\\}
\begin{document}
\cabecalhoInstitucional
% Base centralizada das 60 empresas prospectadas — contato e estratégia.
% Verificação dos contatos: 13/09/2026.
% Regra: dados específicos somente quando confirmados; caso contrário, o registro
% indica canal regional/distribuidor e orienta confirmação antes da expedição.
%
% Cada empresa é um registro de sete campos, um por linha:
%   \ContatoEmpresa
%     {chave}          nome do arquivo do ofício, sem extensão
%     {empresa}        razão/nome usado nos documentos
%     {nível}          P1, P2 ou P3 (o texto do rótulo está em \rotuloPrioridade)
%     {estratégia}     abordagem pretendida na prospecção
%     {situação}       situação do contato no Brasil
%     {setor}          setor prioritário, impresso no campo A/C do ofício
%     {canal}          endereço, telefone, e-mail ou canal oficial no Brasil
%
% Esta é a fonte única desses dados: os arquivos de ofício não os repetem.
% O ofício usa apenas \contatoPrioritario (o setor a quem o documento é
% endereçado). Nível, estratégia, situação e canal são informação interna de
% trabalho e aparecem somente no diretório (empresas/diretorio_contatos.tex).

\providecommand{\setorContato}{Relações Institucionais}
\providecommand{\siteOficial}{canal oficial da empresa}

% Rótulo de cada nível de prioridade, definido em um único lugar.
\newcommand{\rotuloPrioridade}[1]{%
  \ifcsdef{prioridade@rotulo@#1}{\csuse{prioridade@rotulo@#1}}{#1}%
}
\csdef{prioridade@rotulo@P1}{P1 — Máxima}
\csdef{prioridade@rotulo@P2}{P2 — Alta}
\csdef{prioridade@rotulo@P3}{P3 — Oportunidade}

% Valores padrão, usados quando a empresa ainda não possui registro na base.
\newcommand{\statusContatoBrasil}{canal institucional}
\newcommand{\contatoPrioritario}{\setorContato}
\newcommand{\contatoBrasil}{\siteOficial}
\newcommand{\empresaContato}{empresa sem registro na base}
\newcommand{\prioridade}{não classificada}
\newcommand{\estrategia}{estratégia a definir}

\newcommand{\emailContato}[1]{\href{mailto:#1}{\nolinkurl{#1}}}

% Lista de chaves na ordem de cadastro, usada pelo diretório interno.
\newcommand{\listaContatos}{}

% Cadastro de um registro. Os campos ficam guardados em macros indexadas pela
% chave, o que permite tanto a seleção individual quanto a varredura completa.
\newcommand{\ContatoEmpresa}[7]{%
  \listgadd{\listaContatos}{#1}%
  \csdef{contato@#1@empresa}{#2}%
  \csdef{contato@#1@nivel}{#3}%
  \csdef{contato@#1@estrategia}{#4}%
  \csdef{contato@#1@status}{#5}%
  \csdef{contato@#1@setor}{#6}%
  \csdef{contato@#1@canal}{#7}%
}

% Seleção pela chave. Silenciosamente mantém os padrões se a chave não existir.
\newcommand{\selecionaContato}[1]{%
  \ifcsdef{contato@#1@setor}{%
    \renewcommand{\empresaContato}{\csuse{contato@#1@empresa}}%
    \renewcommand{\prioridade}{\rotuloPrioridade{\csuse{contato@#1@nivel}}}%
    \renewcommand{\estrategia}{\csuse{contato@#1@estrategia}}%
    \renewcommand{\statusContatoBrasil}{\csuse{contato@#1@status}}%
    \renewcommand{\contatoPrioritario}{\csuse{contato@#1@setor}}%
    \renewcommand{\contatoBrasil}{\csuse{contato@#1@canal}}%
  }{}%
}

% Seleção automática pelo nome do documento em compilação (\jobname).
% Aceita tanto "01_schneider_electric" quanto "empresas/01_schneider_electric",
% cobrindo o Overleaf e a compilação pela raiz do repositório.
% Observação: \jobname precisa ser expandido antes da comparação; por isso a
% busca é feita via \csname, e não com testes de string sobre \jobname.
\newcommand{\selecionaContatoPeloArquivo}{%
  \selecionaContato{\jobname}%
  \selecionaContato{empresas/\jobname}%
}

\ContatoEmpresa
  {01_schneider_electric}
  {Schneider Electric}
  {P1}
  {Doação direta e parceria educacional para modernização de laboratório.}
  {oficial direto no Brasil}
  {Sustentabilidade / Access to Education / Relações Institucionais}
  {Schneider Electric Brasil — unidade Cajamar/SP: Av. Marginal do Ribeirão dos Cristais, 200, Prédio 400, Jordanésia, Cajamar/SP, CEP 07775-240. Customer Care: \emailContato{ccc.br@schneider-electric.com}; canal oficial: \url{https://www.se.com/br/pt/about-us/sustainability/access-to-education/educational-equipment-services/}.}

\ContatoEmpresa
  {02_vale}
  {Vale}
  {P1}
  {Projeto estruturante de modernização tecnológica dos laboratórios do DEE.}
  {oficial direto no Brasil}
  {Relações Institucionais / Investimento Social / Relacionamento com Comunidades}
  {Vale Brasil — atuação direta no Maranhão. Alô Vale / Relacionamento com Comunidades: 0800 285 7000; Fale Conosco institucional: \url{https://vale.com/pt/fale-conosco}.}

\ContatoEmpresa
  {03_eneva}
  {Eneva}
  {P1}
  {Parceria educacional para implantação de laboratório de energia e geração.}
  {oficial direto no Brasil}
  {Relações Institucionais / ESG / Investimento Social}
  {ENEVA — sede: Praia de Botafogo, 501, Torre Corcovado, sala 404 B, Rio de Janeiro/RJ, CEP 22250-040. Central: 0800 730 1060; contato institucional: \url{https://www.eneva.com.br/contato/}.}

\ContatoEmpresa
  {04_equatorial_energia}
  {Equatorial Energia}
  {P1}
  {Parceria institucional para laboratório de Smart Grid, proteção e medição.}
  {oficial direto no Maranhão}
  {Instituto Equatorial / Relações Institucionais / Projetos Sociais}
  {Equatorial Maranhão — atendimento local: 116; WhatsApp: (98) 2055-0116; clientes corporativos: 0800 280 2800, \emailContato{grandesclientes.maranhao@equatorialenergia.com.br}. Para parceria, solicitar direcionamento ao Instituto Equatorial/Relações Institucionais.}

\ContatoEmpresa
  {05_weg}
  {WEG}
  {P1}
  {Laboratório WEG de Máquinas e Acionamentos Elétricos.}
  {oficial direto no Brasil}
  {Relações Institucionais / Educação / Automação e Drives}
  {WEG — Av. Pref. Waldemar Grubba, 3300, Vila Lalau, Jaraguá do Sul/SC, CEP 89256-900. Tel.: +55 (47) 3276-4000; \emailContato{info-br@weg.net}.}

\ContatoEmpresa
  {06_rockwell_automation}
  {Rockwell Automation}
  {P1}
  {Célula didática Allen-Bradley de automação.}
  {oficial direto no Brasil}
  {Relações Institucionais / Education / Workforce Development}
  {Rockwell Automation Brasil — atendimento de vendas e parceiros pelo canal oficial brasileiro: \url{https://www.rockwellautomation.com/pt-br/company/about-us/contact-us.html}. Escritório em São Paulo/SP; solicitar encaminhamento para Education/Workforce Development.}

\ContatoEmpresa
  {07_siemens}
  {Siemens}
  {P1}
  {Laboratório de Automação e Digitalização.}
  {oficial direto no Brasil}
  {Relações Institucionais / Educação / Digital Industries}
  {Siemens Brasil — Av. Mutinga, 3800, Pirituba, São Paulo/SP, CEP 05110-902. SAC: 0800 11 9484; \emailContato{atendimento.br@siemens.com}.}

\ContatoEmpresa
  {08_altus}
  {Altus Sistemas de Automação}
  {P1}
  {Bancadas nacionais de automação.}
  {oficial direto no Brasil e representante no Maranhão}
  {Relações Institucionais / Educação / Comercial Maranhão}
  {Altus — Av. Theodomiro Porto da Fonseca, 3101, Lote 01, São Leopoldo/RS, CEP 93022-715. Tel.: +55 (51) 3589-9500; \emailContato{vendas@altus.com.br}. Representante Maranhão: Maicon Lima, +55 (51) 99128-4190, \emailContato{maicon.lima@altus.com.br}.}

\ContatoEmpresa
  {09_intelbras}
  {Intelbras}
  {P1}
  {Estruturação de laboratório de redes, IoT e telecom.}
  {oficial direto no Brasil e unidade no Nordeste}
  {Relações Institucionais / Educação / iTEC}
  {Intelbras — matriz em São José/SC, tel. (48) 3281-9500. Unidade Nordeste: Rua Riachão, 200, Módulo 1C, Muribeca, Jaboatão dos Guararapes/PE, CEP 54355-057, tel. (81) 4042-3331. iTEC: (48) 3281-9529.}

\ContatoEmpresa
  {10_petrobras}
  {Petrobras}
  {P1}
  {Monitoramento de chamadas públicas e doação de bens.}
  {programa público oficial no Brasil}
  {Doação de Bens / Relacionamento Institucional}
  {Petrobras — canal oficial de Doação de Bens: \url{https://petrobras.com.br/negocios/doacao-de-bens}. SAC: 0800 728 9001. A habilitação depende de cada chamada pública e respectivo público elegível.}

\ContatoEmpresa
  {11_festo}
  {Festo}
  {P2}
  {Parceria Festo Didactic.}
  {oficial direto no Brasil}
  {Festo Didactic Brasil / Educação}
  {Festo Brasil — Rua Giuseppe Crespi, 76, Jardim Santa Emília, São Paulo/SP, CEP 04183-080. Tel.: +55 (11) 5013-1800; Didactic: \emailContato{didactic.br@festo.com}; vendas: \emailContato{vendas@festo.com}.}

\ContatoEmpresa
  {12_smc}
  {SMC}
  {P2}
  {Parceria educacional para bancadas de automação pneumática.}
  {oficial direto no Brasil e filial no Nordeste}
  {Educação / Relações Institucionais / Automação}
  {SMC Brasil — Av. Piraporinha, 777, Jd. Planalto, São Bernardo do Campo/SP, CEP 09891-001. Tel.: +55 (11) 4082-0600; \emailContato{customerservice.br@smc.com}. Filial Recife/PE: Rua Frei Matias Tevis, 285, sala 1201.}

\ContatoEmpresa
  {13_abb}
  {ABB}
  {P2}
  {Cooperação tecnológica multidisciplinar.}
  {oficial direto no Brasil}
  {Relações Institucionais / Educação / Electrification \& Motion}
  {ABB Brasil — Av. Nicolas Boer, 399, 12º andar, Lapa, São Paulo/SP, CEP 01140-060. Contact Center: 0800 014 9111; \emailContato{contact.center@br.abb.com}.}

\ContatoEmpresa
  {14_alcoa_alumar}
  {Alcoa / Alumar}
  {P2}
  {Investimento social e parceria regional.}
  {oficial direto no Maranhão}
  {Alumar / Relações Institucionais / Instituto Alcoa / Investimento Social}
  {Alcoa/Alumar — operação direta em São Luís/MA. Encaminhar ao Alumar/Relações Institucionais e Instituto Alcoa pelo canal institucional: \url{https://www.alcoa.com/brasil/pt/contact}. Confirmar responsável nominal antes da expedição.}

\ContatoEmpresa
  {15_suzano}
  {Suzano}
  {P2}
  {Investimento social e parceria educacional.}
  {oficial direto no Brasil e atuação no Maranhão}
  {Desenvolvimento Social / Doações e Parcerias Sociais}
  {Suzano — canal de Doações e Parcerias Sociais. Para Maranhão: 0800 771 1418; portal institucional: \url{https://www.suzano.com.br/}.}

\ContatoEmpresa
  {16_hitachi_energy}
  {Hitachi Energy}
  {P2}
  {Laboratório de proteção e automação de subestações.}
  {oficial direto no Brasil}
  {Relações Institucionais / Educação / Grid Automation}
  {Hitachi Energy Brasil — Customer Connect: +55 (11) 4210-4899; consulta comercial oficial: \url{https://www.hitachienergy.com/br/pt/contact-us}. Operações no Brasil em Guarulhos/SP e Blumenau/SC.}

\ContatoEmpresa
  {17_eaton}
  {Eaton}
  {P2}
  {Parceria educacional em proteção e energia.}
  {oficial direto no Brasil}
  {Eaton na Educação / Relações Institucionais / Electrical}
  {Eaton Brasil — atendimento comercial: 0800 00 32866, \emailContato{vendaseaton@eaton.com}; suporte técnico: +55 (15) 3481-9250, \emailContato{CATbrasil@eaton.com}.}

\ContatoEmpresa
  {18_phoenix_contact}
  {Phoenix Contact}
  {P2}
  {Bancada de conectividade e IIoT.}
  {oficial direto no Brasil}
  {Educação / Relações Institucionais / Automação}
  {Phoenix Contact Ind. Com. Ltda. — Av. das Nações Unidas, 11.541, 19º andar, Cidade Monções, São Paulo/SP, CEP 04578-000. Tel.: +55 (11) 3871-6400; \emailContato{vendas@phoenixcontact.com.br}.}

\ContatoEmpresa
  {19_emerson}
  {Emerson}
  {P2}
  {Laboratório de instrumentação industrial.}
  {oficial direto no Brasil}
  {Serviços Educacionais / Relações Institucionais / Automation Solutions}
  {Emerson Brasil — atendimento/suporte no Brasil: +55 (15) 3413-8944 / 8945; contato corporativo e vendas: \url{https://www.emerson.com/pt/corporate/contact-us}. Solicitar encaminhamento a Serviços Educacionais/Automation Solutions.}

\ContatoEmpresa
  {20_fluke}
  {Fluke}
  {P2}
  {Kit institucional de instrumentos de medição.}
  {oficial direto no Brasil}
  {Relações Institucionais / Educação / Vendas Brasil}
  {Fluke do Brasil — CENESP, Av. Maria Coelho Aguiar, 215, Bloco G, 1º andar, São Paulo/SP, CEP 05804-900. Tel.: (11) 3530-8901; \emailContato{insidesales-br@fluke.com}.}

\ContatoEmpresa
  {21_keysight}
  {Keysight Technologies}
  {P3}
  {Programa universitário e apoio a laboratório de eletrônica.}
  {oficial direto no Brasil}
  {Education / University Program / Vendas Brasil}
  {Keysight Brasil — Av. Marcos Penteado de Ulhoa Rodrigues, 939, 6º andar, Sala A, Tamboré, Barueri/SP. Tel.: +55 (11) 3351-7010; \emailContato{lar_orders@keysight.com}.}

\ContatoEmpresa
  {22_tektronix_keithley}
  {Tektronix / Keithley}
  {P3}
  {Apoio ao laboratório de eletrônica e medidas.}
  {parceiro/distribuidor autorizado no Brasil}
  {Education / Academic Program / Parceiro autorizado Brasil}
  {Tektronix/Keithley — selecionar parceiro autorizado para o Brasil no localizador oficial: \url{https://www.tek.com/en/buy/partner-locator}. Encaminhar também cópia ao programa educacional/global da fabricante.}

\ContatoEmpresa
  {23_national_instruments}
  {National Instruments (NI)}
  {P3}
  {Laboratório de aquisição de dados e instrumentação virtual.}
  {canal Brasil integrado à Emerson}
  {Academic / Education / Test \& Measurement}
  {National Instruments (NI), marca da Emerson — atendimento de Teste e Medição pelo canal oficial: \url{https://www.ni.com/pt-br/support/contact-us.html} e \url{https://www.emerson.com/pt/corporate/contact-us}. Solicitar equipe Academic/Education.}

\ContatoEmpresa
  {24_yaskawa}
  {Yaskawa}
  {P3}
  {Célula didática de robótica e acionamentos.}
  {oficial direto no Brasil}
  {Relações Institucionais / Educação / Robótica e Drives}
  {Yaskawa Brasil — Av. Piraporinha, 777, Vila Nogueira, Diadema/SP, CEP 09950-000. Tel.: +55 (11) 3585-1100; \emailContato{comercial@yaskawa.com.br}.}

\ContatoEmpresa
  {25_mitsubishi_electric}
  {Mitsubishi Electric}
  {P3}
  {Célula de automação integrada.}
  {oficial direto no Brasil}
  {Factory Automation / Educação / Relações Institucionais}
  {Mitsubishi Electric do Brasil — Av. Adelino Cardana, 293, 21º andar, Bethaville, Barueri/SP. Tel.: +55 (11) 4689-3000; canal: \url{https://br.mitsubishielectric.com/}.}

\ContatoEmpresa
  {26_bosch_rexroth}
  {Bosch Rexroth}
  {P3}
  {Laboratório de acionamentos e hidráulica.}
  {oficial direto no Brasil}
  {Bosch Rexroth Academy / Didática / Educação}
  {Bosch Rexroth Ltda. — Rod. Dom Pedro I, km 97, Itatiba/SP, CEP 13252-800. Tel.: 0800 704 5446; Didática: \emailContato{didatica@boschrexroth.com.br}; \emailContato{Rexroth.FastChat@boschrexroth.com.br}.}

\ContatoEmpresa
  {27_cisco}
  {Cisco}
  {P3}
  {Laboratório de redes e cibersegurança.}
  {oficial direto no Brasil}
  {Cisco Networking Academy / Educação / Relações Institucionais}
  {Cisco do Brasil — Av. das Nações Unidas, 12.901, Torre Oeste, São Paulo/SP, CEP 04578-903. Vendas Brasil: 0800 891 4972 (opção 1); contato: \url{https://www.cisco.com/c/pt_br/about/contact-cisco.html}.}

\ContatoEmpresa
  {28_huawei}
  {Huawei}
  {P3}
  {ICT Academy e parceria tecnológica.}
  {oficial direto no Brasil e escritório no Nordeste}
  {Huawei ICT Academy / Enterprise / Relações Institucionais}
  {Huawei Brasil — São Paulo: Rua Arquiteto Olavo Redig de Campos, 105, tel. (11) 5105-5105. Escritório Recife: Av. Eng. Domingos Ferreira, 2589, Boa Viagem, Recife/PE. Enterprise: 0800 595 3456; \emailContato{br-support@huawei.com}.}

\ContatoEmpresa
  {29_dell_technologies}
  {Dell Technologies}
  {P3}
  {Modernização computacional dos laboratórios.}
  {oficial direto no Brasil}
  {Educação / Relações Institucionais / Vendas Corporativas}
  {Dell Brasil — Av. Industrial Belgraf, 400, Eldorado do Sul/RS, CEP 92990-000. Vendas: 0800 970 2017; suporte: 0800 722 3300; portal: \url{https://www.dell.com/pt-br}.}

\ContatoEmpresa
  {30_lenovo}
  {Lenovo}
  {P3}
  {Modernização de laboratórios computacionais.}
  {oficial direto no Brasil}
  {Educação / Vendas B2B / Relações Institucionais}
  {Lenovo Brasil — Estr. Municipal José Costa de Mesquita, 200, Mód. 11, Indaiatuba/SP, CEP 13337-200. Vendas B2B: 0800 536 6861 (opção 2); \url{https://www.lenovo.com/br/pt/contact/}.}

\ContatoEmpresa
  {31_fnirsi}
  {FNIRSI}
  {P2}
  {Solicitar um conjunto educacional amplo de instrumentação para os laboratórios de Eletrônica, Medidas Elétricas, Automação, Sistemas Digitais e projetos de pesquisa.}
  {sem filial brasileira oficial confirmada}
  {Business / Bulk Order / Parceria Educacional}
  {FNIRSI — contato direto com a fabricante: \emailContato{business@fnirsi.com}; \url{https://www.fnirsi.com/pages/contact}. Em paralelo, consultar revenda/distribuidor brasileiro vigente antes do envio.}

\ContatoEmpresa
  {32_sel}
  {Schweitzer Engineering Laboratories — SEL}
  {P1}
  {Solicitar formalmente participação no programa de doação de equipamentos universitários da SEL.}
  {oficial direto no Brasil}
  {University Relations / Universidade SEL / Vendas Institucionais}
  {SEL Brasil — Av. Pierre Simon de Laplace, 633, Techno Park, Campinas/SP, CEP 13069-320. Vendas: +55 (19) 3518-2110, \emailContato{vendas@selinc.com}; atendimento: +55 (19) 3515-2060, \emailContato{atendimento@selinc.com}.}

\ContatoEmpresa
  {33_rigol}
  {RIGOL Technologies}
  {P2}
  {Solicitar parceria educacional e eventual doação ou condição institucional para modernização dos laboratórios de Eletrônica e Medidas.}
  {canal oficial específico para o Brasil}
  {Education / University Program / Vendas Brasil}
  {RIGOL — canal Brasil: \emailContato{info.br@rigol.com}; suporte global: \emailContato{service.global@rigol.com}; revendedores oficiais: \url{https://int.rigol.com/CONTACT/where-to-buy.html}.}

\ContatoEmpresa
  {34_siglent}
  {SIGLENT Technologies}
  {P2}
  {Solicitar parceria educacional para equipar bancadas de Eletrônica e Medidas com instrumentos de teste e medição.}
  {site Brasil e distribuição local}
  {Education / Vendas Brasil / Distribuidor autorizado}
  {SIGLENT — portal Brasil: \url{https://www.siglent.com/br/}. Encaminhar pelo contato comercial/distribuidor brasileiro indicado no site e copiar a área educacional/global da fabricante. Confirmar representante vigente antes da expedição.}

\ContatoEmpresa
  {35_endress_hauser}
  {Endress+Hauser}
  {P1}
  {Solicitar parceria educacional e doação de instrumentos, demonstradores ou bancada didática para Instrumentação e Automação.}
  {oficial direto no Brasil}
  {Educação / Relações Institucionais / Training}
  {Endress+Hauser Brasil — Estr. Mun. Antônio Sesti, 600, Itatiba/SP, CEP 13254-085. Tel.: +55 (11) 5033-4333; \emailContato{info.br@endress.com}.}

\ContatoEmpresa
  {36_universal_robots}
  {Universal Robots}
  {P2}
  {Solicitar parceria educacional, kit de ensino e capacitação docente para implantação de célula de robótica colaborativa.}
  {oficial direto no Brasil}
  {Education / Relações Institucionais / Vendas Brasil}
  {Universal Robots Brasil — Rua Alegre, 470, cj. 1102, São Caetano do Sul/SP, CEP 09550-250. Tel.: 0800 777 0442; \emailContato{sales@universal-robots.com}.}

\ContatoEmpresa
  {37_fanuc}
  {FANUC}
  {P2}
  {Solicitar parceria educacional para implantação de célula robótica didática e acesso ao programa de formação/certificação.}
  {oficial direto no Brasil}
  {Education / Relações Institucionais / Robótica}
  {FANUC South America — Av. Emb. Macedo Soares, 10001, Prédio 6, São Paulo/SP, CEP 05095-035. Canal oficial: \url{https://www.fanucamerica.com/contact-us}; solicitar equipe Education/Robotics Brasil.}

\ContatoEmpresa
  {38_kuka}
  {KUKA}
  {P2}
  {Solicitar parceria educacional para célula robótica didática e capacitação em automação industrial.}
  {oficial direto no Brasil}
  {KUKA Education / Relações Institucionais / Vendas}
  {KUKA Roboter do Brasil — Av. Osvaldo Fregonezi, 171, São Bernardo do Campo/SP, CEP 09851-015. Vendas: +55 (11) 4942-8299, \emailContato{vendas.br@kuka.com}.}

\ContatoEmpresa
  {39_omron}
  {OMRON Automation}
  {P2}
  {Prospecção estratégica para modernização do Laboratório de Automação, sem pressupor programa público permanente de doação.}
  {oficial direto no Brasil}
  {Educação / Automation Center / Relações Institucionais}
  {OMRON Eletrônica do Brasil — Alameda Vicente Pinzón, 51, 3º andar, Vila Olímpia, São Paulo/SP, CEP 04547-130. Tel.: +55 (11) 5171-8920; Automation/PoC São Paulo: \url{https://automation.omron.com/}.}

\ContatoEmpresa
  {40_beckhoff}
  {Beckhoff Automation}
  {P2}
  {Prospecção de plataforma didática de automação baseada em PC e EtherCAT.}
  {oficial direto no Brasil}
  {Educação / Relações Institucionais / Automação}
  {Beckhoff Automação Industrial Ltda. — Rua Caminho do Pilar, 1362, Vila Gilda, Santo André/SP, CEP 09190-000. Tel.: +55 (11) 4126-3232; \emailContato{info@beckhoff.com.br}.}

\ContatoEmpresa
  {41_wago}
  {WAGO}
  {P2}
  {Prospecção estratégica para bancadas de automação, conectividade e redes industriais.}
  {oficial direto no Brasil}
  {Educação / Relações Institucionais / Vendas}
  {WAGO Eletroeletrônicos Ltda. — Rua Trípoli, 640, Jundiaí/SP, CEP 13212-217. Tel.: +55 (11) 2923-7200; \emailContato{vendas.br@wago.com}; \emailContato{info.br@wago.com}.}

\ContatoEmpresa
  {42_sick}
  {SICK}
  {P2}
  {Prospecção de kits didáticos de sensoriamento, segurança e identificação industrial.}
  {oficial direto no Brasil}
  {Educação / Sensores / Relações Institucionais}
  {SICK Solução em Sensores Ltda. — Av. dos Imarés, 391, Indianópolis, São Paulo/SP, CEP 04085-000. Tel.: +55 (11) 3215-4900; \emailContato{comercial@sick.com.br}.}

\ContatoEmpresa
  {43_ifm}
  {ifm electronic}
  {P2}
  {Prospecção estratégica para laboratório de sensores, IO-Link e manutenção preditiva.}
  {oficial direto no Brasil e escritório em Fortaleza}
  {Educação / Escritório de Vendas Fortaleza / Relações Institucionais}
  {ifm electronic — escritório Fortaleza: Av. Washington Soares, 3663, WSTC Torre II, sala 1404, Edson Queiroz, Fortaleza/CE, CEP 60811-341. Tel.: 0800 544 2436; \emailContato{info.br@ifm.com}.}

\ContatoEmpresa
  {44_balluff}
  {Balluff}
  {P2}
  {Prospecção de kits educacionais de sensoriamento, IO-Link e identificação industrial.}
  {oficial direto no Brasil}
  {Educação / IO-Link / Relações Institucionais}
  {Balluff Controles Elétricos Ltda. — Rod. Eng. Ermênio de Oliveira Penteado, km 57,5, Indaiatuba/SP, CEP 13337-300. Tel.: +55 (19) 3876-9999; contato comercial via \url{https://www.balluff.com/pt-br/contact}.}

\ContatoEmpresa
  {45_pepperl_fuchs}
  {Pepperl+Fuchs}
  {P2}
  {Prospecção estratégica de kits de sensores e interfaces para Automação e Instrumentação.}
  {oficial direto no Brasil}
  {Educação / Sensores e Automação de Processo / Vendas}
  {Pepperl+Fuchs Brasil — atendimento comercial: +55 (11) 4007-1448; \emailContato{vendas@br.pepperl-fuchs.com}; portal: \url{https://www.pepperl-fuchs.com/brazil/}.}

\ContatoEmpresa
  {46_weidmueller}
  {Weidmüller}
  {P2}
  {Prospecção de componentes e controladores para painéis didáticos e redes industriais.}
  {oficial direto no Brasil}
  {Educação / Relações Institucionais / Vendas}
  {Weidmüller Conexel do Brasil — Av. Presidente Juscelino, 642, Diadema/SP, CEP 09950-370. Central: +55 (11) 4366-9600; vendas: \emailContato{vendas@weidmueller.com}.}

\ContatoEmpresa
  {47_delta_electronics}
  {Delta Electronics}
  {P2}
  {Prospecção de plataforma integrada de CLP, HMI, inversor e servo para laboratório.}
  {oficial direto no Brasil}
  {Educação / Automação Industrial / Relações Institucionais}
  {Delta Electronics Brasil — Estrada Velha Rio-São Paulo, 5300, Eugênio de Melo, São José dos Campos/SP, CEP 12247-001. Tel.: +55 (12) 3932-2300; \emailContato{contato.automacao@deltaww.com}.}

\ContatoEmpresa
  {48_yokogawa}
  {Yokogawa}
  {P2}
  {Prospecção de instrumentação de alta precisão para Medidas, Energia e Automação.}
  {oficial direto no Brasil}
  {Educação / Instrumentação / Relações Institucionais}
  {Yokogawa América do Sul — Alameda Xingu, 850, Alphaville, Barueri/SP, CEP 06455-030. Tel.: +55 (11) 5681-2400 / 3513-1300; \emailContato{YSA-eproc@yokogawa.com}.}

\ContatoEmpresa
  {49_hioki}
  {HIOKI}
  {P2}
  {Prospecção de instrumentação para potência, manutenção e aquisição de dados.}
  {atendimento regional oficial para o Brasil}
  {Education / Sales / Distribuidor Brasil}
  {HIOKI — o Brasil é atendido oficialmente pela HIOKI USA Corporation para América Latina; tel. +1 214 281 4546 e formulário em \url{https://www.hioki.com/br-pt/support/network}. Usar também distribuidor autorizado brasileiro indicado no locator vigente.}

\ContatoEmpresa
  {50_megger}
  {Megger}
  {P2}
  {Prospecção de instrumentos para ensaios de instalações, máquinas, transformadores e proteção.}
  {oficial direto no Brasil}
  {Educação / Vendas Brasil / Sistemas de Potência}
  {Megger Brazil — Rua Bom Sucesso, 220, sala 3801, Tatuapé, São Paulo/SP, CEP 03305-000. Tel.: +55 (11) 96335-3518; contato oficial: \url{https://www.megger.com/}.}

\ContatoEmpresa
  {51_minipa}
  {Minipa}
  {P2}
  {Prospecção nacional para kits de instrumentos de bancada e portáteis destinados ao ensino técnico e superior.}
  {oficial direto no Brasil}
  {Educação / Relações Institucionais / Vendas}
  {Minipa do Brasil — São Paulo/SP. Tel.: (11) 5078-1850; \emailContato{minipa@minipa.com.br}; \url{https://www.minipa.com.br/}.}

\ContatoEmpresa
  {52_instrutherm}
  {Instrutherm}
  {P3}
  {Prospecção complementar para instrumentos portáteis de ensino, manutenção e diagnóstico.}
  {oficial direto no Brasil}
  {Relações Institucionais / Licitações / Vendas}
  {Instrutherm — Rua Jorge de Freitas, 274, Freguesia do Ó, São Paulo/SP, CEP 02911-030. Central: (11) 2144-2800; Licitação: (11) 2144-2849; contato: \url{https://www.instrutherm.com.br/contato}.}

\ContatoEmpresa
  {53_pico_technology}
  {Pico Technology}
  {P3}
  {Prospecção para laboratórios flexíveis de aquisição de dados e osciloscopia baseada em computador.}
  {distribuidor/representante para o Brasil}
  {Education / Distribuidor local / Vendas}
  {Pico Technology — localizar distribuidor vigente para o Brasil em \url{https://www.picotech.com/distributors}. Encaminhar cópia ao fabricante/education para avaliação institucional.}

\ContatoEmpresa
  {54_rohde_schwarz}
  {Rohde \& Schwarz}
  {P1}
  {Prospecção educacional prioritária, aproveitando o portfólio formal da empresa para laboratórios de ensino e pesquisa.}
  {oficial direto no Brasil}
  {University Program / Education / Vendas Brasil}
  {Rohde \& Schwarz do Brasil — Av. Magalhães de Castro, 4800, Continental Tower, 17º andar, São Paulo/SP, CEP 05676-120. Vendas: +55 (11) 2246-0076; \emailContato{vendas.brasil@rohde-schwarz.com}.}

\ContatoEmpresa
  {55_anritsu}
  {Anritsu}
  {P2}
  {Prospecção estratégica para fortalecer Telecomunicações, RF e atividades práticas de análise de espectro e redes.}
  {oficial direto no Brasil}
  {Education / Relações Institucionais / Vendas Brasil}
  {Anritsu Eletrônica Ltda. — Praça Amadeu Amaral, 27, 1º andar, Bela Vista, São Paulo/SP, CEP 01327-010. Tel.: +55 (11) 3283-2511; \url{https://www.anritsu.com/pt-BR}.}

\ContatoEmpresa
  {56_viavi}
  {VIAVI Solutions}
  {P2}
  {Prospecção estratégica para laboratórios de Telecomunicações, redes ópticas e comunicações móveis.}
  {oficial direto no Brasil}
  {Education / Vendas Brasil / Telecomunicações}
  {VIAVI Solutions do Brasil — Av. Eng. Luís Carlos Berrini, 936, 9º andar, São Paulo/SP, CEP 04571-000. Tel.: +55 (11) 5503-3800; \emailContato{vendas.brasil@viavisolutions.com}.}

\ContatoEmpresa
  {57_infineon}
  {Infineon Technologies}
  {P1}
  {Prospecção educacional prioritária por meio da University Alliance, kits acadêmicos, amostras e cooperação em Eletrônica de Potência.}
  {canal acadêmico global com presença comercial no Brasil}
  {University Alliance / Academic Program / South America}
  {Infineon — priorizar o University Alliance para kits, materiais e amostras: \url{https://www.infineon.com/university}. Há presença comercial em São Paulo; confirmar endereço/telefone vigentes no momento do envio pelo contato regional.}

\ContatoEmpresa
  {58_stmicroelectronics}
  {STMicroelectronics}
  {P2}
  {Prospecção acadêmica para kits, placas de avaliação e cooperação em Eletrônica de Potência e sistemas embarcados.}
  {presença comercial em São Paulo; contato local a confirmar}
  {University Program / Power Electronics / Sales Brazil}
  {STMicroelectronics — equipe de aplicação e vendas com presença em São Paulo/SP. Encaminhar pelo canal oficial \url{https://www.st.com/content/st_com/en/contact-us.html} e solicitar University Program/Power Conversion Brasil; confirmar dados locais antes da expedição.}

\ContatoEmpresa
  {59_texas_instruments}
  {Texas Instruments}
  {P1}
  {Prospecção educacional prioritária com base no portfólio universitário de Eletrônica de Potência e plataformas de controle digital.}
  {canal acadêmico global; atendimento ao Brasil}
  {University Program / Power Electronics / Education}
  {Texas Instruments — priorizar TI University Program e Power Management Lab Kit: \url{https://www.ti.com/university}. Para contato comercial, usar \url{https://www.ti.com/info/contact-us.html}; confirmar responsável Brasil antes da expedição.}

\ContatoEmpresa
  {60_onsemi}
  {onsemi}
  {P2}
  {Prospecção educacional e de apoio a laboratório, considerando iniciativas STEAM e cooperações universitárias da empresa.}
  {atendimento regional/distribuidores; filial brasileira não confirmada}
  {STEAM / Community Giving / University Relations}
  {onsemi — Sales Offices e distribuidores autorizados: \url{https://www.onsemi.com/support/sales}. Para apoio educacional, encaminhar também ao programa STEAM/Community Giving; não foi confirmada filial brasileira própria nesta verificação.}



\noindent{\sffamily\bfseries OFÍCIO Nº \campo{número/ano} — DEE/MTC/IFMA}

\vspace{1.5mm}
\begin{flushright}
São Luís — MA, \campo{data}.
\end{flushright}

\vspace{4mm}
\noindent À\\
\textbf{\empresa}\\
\textbf{A/C:} \contatoPrioritario

\vspace{5mm}
\noindent\textbf{Assunto:} Solicitação de apoio para modernização dos laboratórios do Departamento de Eletroeletrônica do IFMA.

\vspace{4mm}
Prezados(as) Senhores(as),

O Departamento de Eletroeletrônica -- DEE do IFMA -- Campus São Luís/Monte Castelo vem apresentar seu Programa de Modernização Tecnológica dos Laboratórios e consultar a \textbf{\empresa} sobre a possibilidade de colaboração com equipamentos, kits didáticos, softwares, licenças, treinamento ou outra forma de cooperação compatível com sua atuação.

O DEE atende cursos técnicos de Eletrotécnica, Eletrônica e Automação e o curso de Engenharia Elétrica. Seus laboratórios são utilizados em aulas práticas, iniciação científica, TCCs, projetos de pesquisa e extensão e também podem apoiar trabalhos vinculados a mestrado e doutorado.

A modernização permitirá que os estudantes tenham contato com tecnologias próximas das utilizadas no setor produtivo, fortalecendo a formação de profissionais capazes de configurar, operar, ensaiar, medir, diagnosticar e utilizar esses equipamentos e sistemas com segurança e domínio técnico.

Pela atuação da \textbf{\empresa} em \areaEmpresa, entendemos que há uma aproximação natural entre suas soluções e algumas das necessidades do Departamento. \textoCooperacao

\vspace{2mm}
\noindent{\sffamily\bfseries\color{cinzaescuro} Equipamentos e soluções de interesse}
\vspace{1.5mm}

\renewcommand{\arraystretch}{1.18}
\begin{tabularx}{\textwidth}{>{\RaggedRight\arraybackslash}X >{\centering\arraybackslash}p{1.4cm} >{\RaggedRight\arraybackslash}X}
\toprule
\textbf{Equipamento/solução} & \textbf{Qtd.} & \textbf{Aplicação acadêmica}\\
\midrule
\itensEmpresa
\bottomrule
\end{tabularx}

\vspace{4mm}
A relação pode ser ajustada de acordo com a disponibilidade e o portfólio da empresa. Também temos interesse em unidades de demonstração, equipamentos de treinamento, itens descontinuados em boas condições de funcionamento e outros recursos que possam ser utilizados com segurança em atividades acadêmicas.

Além do uso em ensino, os recursos poderão apoiar pesquisa aplicada, dissertações, teses e projetos de \textbf{Pesquisa e Desenvolvimento (P\&D)}. O Departamento está à disposição para discutir projetos conjuntos, provas de conceito, protótipos, ensaios, medições e validações em temas de interesse comum.

As doações e parcerias serão formalizadas pelos setores competentes do IFMA, conforme os procedimentos institucionais aplicáveis.

Ficamos à disposição para apresentar os laboratórios, detalhar as necessidades e conversar sobre possibilidades de cooperação.

\vspace{7mm}
Atenciosamente,

\vspace{13mm}
\begin{center}
\textbf{Francisco dos Santos Viana}\\[-0.2mm]
{\small Chefe do Departamento de Eletroeletrônica — DEE}\\[-0.2mm]
{\small IFMA — Campus São Luís/Monte Castelo}\\[1mm]
{\footnotesize \href{mailto:francisco.viana@ifma.edu.br}{francisco.viana@ifma.edu.br} \quad | \quad (98) 98816-6263}
\end{center}
\end{document}

